% Autómata finito de identificadores y números, y por qué reconoce un lenguaje regular.
% Responsable: Persona 1.
%
% Compilar (dos pasadas, para las referencias cruzadas):
%   pdflatex automata.tex && pdflatex automata.tex
% o bien:
%   latexmk -pdf automata.tex
\documentclass[11pt, a4paper]{article}

\usepackage[T1]{fontenc}
\usepackage[spanish, es-noshorthands, es-tabla]{babel}
\usepackage{lmodern}
\usepackage{microtype}
\usepackage{upquote}
\usepackage[margin=2.3cm]{geometry}
\usepackage{amsmath, amssymb, amsthm}
\usepackage{booktabs}
\usepackage{array}
\usepackage{xcolor}
\usepackage{tikz}
\usetikzlibrary{automata, arrows.meta, positioning}
\usepackage[hidelinks]{hyperref}

\theoremstyle{plain}
\newtheorem{teorema}{Teorema}
\newtheorem{corolario}[teorema]{Corolario}
\theoremstyle{definition}
\newtheorem{definicion}[teorema]{Definición}

\newcommand{\Sig}{\Sigma}
\newcommand{\cls}[1]{\textsf{#1}}            % clase de caracteres
\newcommand{\pto}{\texttt{'.'}}               % el carácter punto
\newcommand{\lex}[1]{\texttt{#1}}             % un lexema
\newcommand{\Lang}{\mathcal{L}}

\title{Autómata finito del analizador léxico:\\identificadores y números}
\author{Compilador Fase 1 --- Persona 1}
\date{Noviembre de 2026}

\begin{document}
\maketitle

\begin{abstract}
Se define formalmente el autómata finito determinista (AFD) que reconoce los
identificadores y los literales numéricos del lenguaje, incluidos los estados de error
del analizador léxico. A partir del teorema de Kleene se demuestra que cada clase de
token es un lenguaje regular, y se verifica en sentido inverso derivando, desde el
autómata, las mismas expresiones regulares de la especificación de tokens.
\end{abstract}

% ======================================================================
\section{Clases de caracteres}

Para que el diagrama sea legible, el alfabeto $\Sig$ (definido en
\nolinkurl{docs/especificacion_tokens.md}) se agrupa en clases. Cada transición etiquetada
con una clase representa una transición por cada carácter de esa clase.

\begin{center}
\begin{tabular}{>{$}c<{$} l l}
  \toprule
  \text{Clase} & Caracteres & Expresión regular \\
  \midrule
  L      & letras ASCII y guion bajo & \texttt{[a-zA-Z\_]} \\
  D      & dígitos                   & \texttt{[0-9]} \\
  \pto   & el punto                  & \texttt{\textbackslash.} \\
  O      & cualquier otro carácter de $\Sig$ (espacios, operadores, delimitadores) & --- \\
  \bottomrule
\end{tabular}
\end{center}

% ======================================================================
\section{Definición formal}

\begin{definicion}[Autómata finito determinista]
Un AFD es una 5-tupla $M = (Q, \Sig, \delta, q_0, F)$ donde $Q$ es un conjunto finito de
estados, $\Sig$ es el alfabeto, $\delta : Q \times \Sig \to Q$ es la función de
transición, $q_0 \in Q$ es el estado inicial y $F \subseteq Q$ es el conjunto de estados
de aceptación. $M$ acepta una cadena $w$ si, partiendo de $q_0$ y leyendo $w$ símbolo a
símbolo, termina en un estado de $F$. El lenguaje de $M$, $\Lang(M)$, es el conjunto de
cadenas que acepta~\cite{sipser, hmu}.
\end{definicion}

El autómata del analizador léxico es
\[
  M = (Q, \Sig, \delta, q_0, F), \qquad
  Q = \{q_0, q_1, q_2, q_3, q_4, q_5, e_1, e_2\}, \qquad
  F = \{q_1, q_2, q_4\},
\]
con $\delta$ dada por la figura~\ref{fig:afd} y la tabla~\ref{tab:delta}.

\begin{figure}[ht]
  \centering
  \resizebox{\linewidth}{!}{% Diagrama del AFD de identificadores y números.
% Lo usan automata.tex (documento) y automata_figura.tex (imagen para el video).
% Requiere: \usetikzlibrary{automata, arrows.meta, positioning}
%
% Clases de caracteres:  L = [a-zA-Z_]   D = [0-9]   '.' = punto
% (dagger) = solo si el punto va seguido de un dígito (un carácter de anticipación)
\begin{tikzpicture}[
    ->, >={Stealth[length=2.6mm]}, shorten >=1pt, auto, semithick,
    every state/.style={minimum size=11mm, inner sep=1pt, font=\normalsize},
    ok/.style={state, accepting, fill=green!10},
    neutral/.style={state, fill=gray!8},
    error/.style={state, draw=red!70!black, fill=red!10, dashed},
    lbl/.style={font=\scriptsize\sffamily, align=center},
    every loop/.style={looseness=6},
  ]
  % Estados
  \node[neutral, initial, initial text={}] (q0) at (0, 0)    {$q_0$};
  \node[ok]                                (q1) at (3, 2.4)  {$q_1$};
  \node[ok]                                (q2) at (3, 0)    {$q_2$};
  \node[neutral]                           (q3) at (5.6, 0)  {$q_3$};
  \node[ok]                                (q4) at (8.2, 0)  {$q_4$};
  \node[neutral]                           (q5) at (10.8, 0) {$q_5$};
  \node[error]                             (e1) at (3, -2.6) {$e_1$};
  \node[error]                             (e2) at (9.5, -2.6) {$e_2$};

  % Etiquetas de lo que reconoce cada estado
  \node[lbl, right=2mm of q1]                {ID / palabra\\reservada};
  \node[lbl, below left=0mm and -1mm of q2]  {NUM\_ENTERO};
  \node[lbl, below=1mm of q3]                {``\texttt{12.}''};
  \node[lbl, below left=0mm and -1mm of q4]  {NUM\_REAL};
  \node[lbl, below=1mm of q5]                {``\texttt{3.5.}''};
  \node[lbl, left=2mm of e1, text=red!60!black]  {error:\\ID empieza\\con dígito};
  \node[lbl, right=2mm of e2, text=red!60!black] {error:\\número mal\\formado};

  % Transiciones
  \path
    (q0) edge              node        {$L$}        (q1)
    (q0) edge              node        {$D$}        (q2)
    (q1) edge [loop above] node        {$L, D$}     (q1)
    (q2) edge [loop above] node        {$D$}        (q2)
    (q2) edge              node        {\texttt{'.'}} (q3)
    (q2) edge              node        {$L$}        (e1)
    (q3) edge              node        {$D$}        (q4)
    (q4) edge [loop above] node        {$D$}        (q4)
    (q4) edge              node        {\texttt{'.'}} (q5)
    (q4) edge              node[swap]  {$L$}        (e2)
    (q5) edge              node        {$D$}        (e2)
    (e1) edge [loop below] node        {$L, D$}     (e1)
    (e1) edge              node[swap]  {\texttt{'.'}$^{\dagger}$} (e2)
    (e2) edge [loop below] node        {$L, D,$ \texttt{'.'}$^{\dagger}$} (e2);
\end{tikzpicture}
}
  \caption{AFD de identificadores y números. Doble círculo: estado de aceptación.
    Línea discontinua: estado de error. $^{\dagger}$: solo si el punto va seguido de un
    dígito (ver la sección~\ref{sec:anticipacion}).}
  \label{fig:afd}
\end{figure}

\begin{table}[htbp]
  \centering
  \caption{Función de transición $\delta$. Una celda vacía (---) significa que no hay
    transición: el lexema termina y se ejecuta la acción de la última columna.}
  \label{tab:delta}
  \begin{tabular}{c c c c c l}
    \toprule
    Estado & $L$ & $D$ & \pto & $O$ & Acción al terminar el lexema \\
    \midrule
    $\to q_0$  & $q_1$ & $q_2$ & ---              & --- & (no es identificador ni número) \\
    $*\,q_1$   & $q_1$ & $q_1$ & ---              & --- & acepta \textsf{ID} o palabra reservada \\
    $*\,q_2$   & $e_1$ & $q_2$ & $q_3$            & --- & acepta \textsf{NUM\_ENTERO} \\
    $q_3$      & ---   & $q_4$ & ---              & --- & error: número mal formado (\lex{12.}) \\
    $*\,q_4$   & $e_2$ & $q_4$ & $q_5$            & --- & acepta \textsf{NUM\_REAL} \\
    $q_5$      & ---   & $e_2$ & ---              & --- & retrocede a $q_4$: \textsf{NUM\_REAL} \\
    $e_1$      & $e_1$ & $e_1$ & $e_2^{\dagger}$  & --- & error: identificador empieza con dígito \\
    $e_2$      & $e_2$ & $e_2$ & $e_2^{\dagger}$  & --- & error: número mal formado \\
    \bottomrule
  \end{tabular}

  \smallskip
  {\footnotesize $\to$ estado inicial \quad $*$ estado de aceptación}
\end{table}

\paragraph{Función de transición parcial.} Las celdas vacías hacen que $\delta$ sea
\emph{parcial}. Para tener un AFD completo en el sentido estricto basta agregar un
estado sumidero $q_\varnothing \notin F$ al que van todas las celdas vacías y que se
transita a sí mismo con todo $\Sig$. Ese estado no cambia el lenguaje aceptado, porque
desde él no se alcanza ningún estado de aceptación~\cite{hmu}.

\paragraph{Estados de error.} $e_1$ y $e_2$ no pertenecen a $F$, así que no cambian
$\Lang(M)$. Su función es práctica: consumen el lexema mal formado completo para
reportar \emph{un solo} error (\lex{3abc} es un error, y no \lex{3} seguido de
\lex{abc}).

\subsection{Anticipación de un carácter}
\label{sec:anticipacion}

El estado $q_5$ y las transiciones marcadas con $^{\dagger}$ modelan la técnica habitual
de los analizadores léxicos: leer hasta donde se pueda y, si el lexema no puede
continuar, \emph{retroceder} al último estado de aceptación visitado~\cite{dragon}. En
\lex{3.5.x}, el autómata pasa por $q_4$, lee el segundo punto y llega a $q_5$. Como
\lex{x} no es un dígito, retrocede a $q_4$: emite \textsf{NUM\_REAL} \lex{3.5} y vuelve a
analizar desde el punto. En el código, esto es la comprobación
\verb|self.current() == "." and is_digit(self.peek())|.

La única excepción deliberada es $q_3$. En \lex{12.} se podría retroceder a $q_2$ y emitir
\lex{12}, pero el equipo decidió reportarlo como número mal formado porque casi siempre
es un número incompleto.

% ======================================================================
\section{Trazas de ejemplo}

\begin{center}
\begin{tabular}{l l l}
  \toprule
  Entrada & Recorrido & Resultado \\
  \midrule
  \lex{x\_1}    & $q_0 \xrightarrow{x} q_1 \xrightarrow{\texttt{\_}} q_1 \xrightarrow{1} q_1$ & \textsf{ID} \\
  \lex{si}      & $q_0 \xrightarrow{s} q_1 \xrightarrow{i} q_1$ & \textsf{ID} $\to$ palabra reservada \textsf{IF} \\
  \lex{42}      & $q_0 \xrightarrow{4} q_2 \xrightarrow{2} q_2$ & \textsf{NUM\_ENTERO} \\
  \lex{3.14}    & $q_0 \xrightarrow{3} q_2 \xrightarrow{\texttt{.}} q_3 \xrightarrow{1} q_4 \xrightarrow{4} q_4$ & \textsf{NUM\_REAL} \\
  \lex{12.}     & $q_0 \xrightarrow{1} q_2 \xrightarrow{2} q_2 \xrightarrow{\texttt{.}} q_3$ & error: termina en $q_3 \notin F$ \\
  \lex{3abc}    & $q_0 \xrightarrow{3} q_2 \xrightarrow{a} e_1 \xrightarrow{b} e_1 \xrightarrow{c} e_1$ & error: ID empieza con dígito \\
  \lex{12.3.4}  & $q_0 \to q_2 \to q_2 \xrightarrow{\texttt{.}} q_3 \xrightarrow{3} q_4 \xrightarrow{\texttt{.}} q_5 \xrightarrow{4} e_2$ & error: número mal formado \\
  \bottomrule
\end{tabular}
\end{center}

% ======================================================================
\section{Por qué el autómata define un lenguaje regular}

\begin{definicion}[Lenguaje regular]
Un lenguaje es \emph{regular} si existe un autómata finito que lo reconoce~\cite{sipser}.
\end{definicion}

\begin{teorema}[Kleene, 1956~\cite{kleene}]
Un lenguaje es aceptado por un autómata finito si y solo si puede describirse con una
expresión regular.
\end{teorema}

Kleene lo enunció en 1956 para los \emph{eventos regulares} de las redes neuronales de
McCulloch y Pitts~\cite{kleene}. La forma actual del teorema viene de tres resultados:

\begin{itemize}
  \item \textbf{Rabin y Scott~\cite{rabinscott}:} los autómatas no deterministas no
    reconocen más lenguajes que los deterministas (construcción de subconjuntos).
  \item \textbf{McNaughton y Yamada~\cite{mcnaughton}:} algoritmos explícitos para
    convertir entre expresiones regulares y diagramas de estados.
  \item \textbf{Thompson~\cite{thompson}:} una construcción de expresión regular a
    autómata que se sigue usando en los generadores de analizadores léxicos.
\end{itemize}

Los libros de texto lo presentan como la equivalencia entre AFD, AFN y expresiones
regulares: \cite[teorema~1.54]{sipser}, \cite[sección~3.2]{hmu} y \cite[secciones~3.6--3.7]{dragon}.

\subsection{Aplicación a nuestros tokens}

Para cada estado de aceptación $q \in F$ sea $\Lang_q(M)$ el conjunto de cadenas que
llevan a $M$ de $q_0$ a $q$. Cada $\Lang_q(M)$ es el lenguaje de un autómata finito: el
mismo $M$ con $F = \{q\}$.

\begin{corolario}
$\Lang_{q_1}(M)$, $\Lang_{q_2}(M)$ y $\Lang_{q_4}(M)$ son lenguajes regulares.
\end{corolario}

\begin{proof}
Cada uno es reconocido por un autómata finito ($M$ tiene 8 estados). Por la definición de
lenguaje regular, los tres lo son. El teorema de Kleene garantiza además que cada uno
tiene una expresión regular equivalente.
\end{proof}

\paragraph{Verificación: del autómata a la expresión regular.} El teorema dice que la
expresión regular existe. Para nuestro autómata se puede obtener a mano eliminando
estados~\cite[sección~3.2.2]{hmu}. Hay que seguir los caminos de $q_0$ a cada estado de
aceptación, donde un lazo se convierte en una estrella de Kleene:

\begin{align*}
  \Lang_{q_1}(M) &= L\,(L \mid D)^{*}
     && \text{($q_0 \xrightarrow{L} q_1$ y el lazo $L, D$ en $q_1$)} \\
  \Lang_{q_2}(M) &= D\,D^{*} = D^{+}
     && \text{($q_0 \xrightarrow{D} q_2$ y el lazo $D$ en $q_2$)} \\
  \Lang_{q_4}(M) &= D^{+}\,\pto\,D\,D^{*} = D^{+}\,\pto\,D^{+}
     && \text{(el camino $q_0 \to q_2 \to q_3 \to q_4$ con sus lazos)}
\end{align*}

Son exactamente las expresiones regulares de \textsf{ID}, \textsf{NUM\_ENTERO} y
\textsf{NUM\_REAL} en \nolinkurl{docs/especificacion_tokens.md}. En nuestro caso, las dos
direcciones del teorema (\emph{autómata $\Rightarrow$ expresión regular} y \emph{expresión
regular $\Rightarrow$ autómata}) llevan al mismo resultado. La prueba
\nolinkurl{tests/test_token_spec.py} verifica lo mismo de forma ejecutable: compara el
analizador léxico, que implementa el autómata, con un tokenizador hecho solo con las
expresiones regulares (\texttt{TOKEN\_PATTERNS} en \nolinkurl{compiler/tokens.py}).

\subsection{Palabras reservadas y el resto de los tokens}

\begin{itemize}
  \item \textbf{Palabras reservadas.} El conjunto $K = \{\lex{entero}, \lex{real},
    \dots, \lex{falso}\}$ tiene 9 cadenas. Todo lenguaje finito es regular. Los
    lenguajes regulares son cerrados bajo complemento e intersección, y por lo tanto bajo
    diferencia~\cite[sección~4.2]{hmu}. Entonces $\Lang(\textsf{ID}) = \Lang_{q_1}(M)
    \setminus K$ también es regular. En el código, esta diferencia es la búsqueda en la
    tabla \texttt{KEYWORDS} después de reconocer el lexema en $q_1$.
  \item \textbf{Operadores, delimitadores, espacios y comentarios.} Cada uno se describe
    con una expresión regular, así que también son regulares.
  \item \textbf{Todos los tokens juntos.} Los lenguajes regulares son cerrados bajo
    unión, de modo que el lenguaje de todos los tokens es regular. Por eso un único
    autómata finito (el analizador léxico) basta para reconocerlos.
\end{itemize}

\subsection{Lo que un autómata finito no puede hacer}

El autómata tiene un número fijo de estados, es decir, \emph{memoria finita}: no puede
contar sin límite. Por eso no puede verificar que los paréntesis $\lex{(}^n\lex{)}^n$
estén balanceados para todo $n$, ni que una variable usada haya sido declarada antes.
El primer problema requiere una gramática libre de contexto (tipo~2) y el segundo va más
allá de ellas~\cite{chomsky56, chomsky59}. Esa es la jerarquía que sigue el compilador:
el analizador léxico es de tipo~3, la sintaxis es de tipo~2 y el análisis semántico
necesita una tabla de símbolos (ver \nolinkurl{docs/reflexion.md}).

% ======================================================================
\begin{thebibliography}{9}

\bibitem{kleene}
S.~C. Kleene.
\newblock Representation of events in nerve nets and finite automata.
\newblock En C.~E. Shannon y J.~McCarthy (eds.), \emph{Automata Studies}, Annals of
  Mathematics Studies~34, pp.~3--41. Princeton University Press, 1956.

\bibitem{rabinscott}
M.~O. Rabin y D.~Scott.
\newblock Finite automata and their decision problems.
\newblock \emph{IBM Journal of Research and Development}, 3(2):114--125, 1959.

\bibitem{mcnaughton}
R.~McNaughton y H.~Yamada.
\newblock Regular expressions and state graphs for automata.
\newblock \emph{IRE Transactions on Electronic Computers}, EC-9(1):39--47, 1960.

\bibitem{thompson}
K.~Thompson.
\newblock Regular expression search algorithm.
\newblock \emph{Communications of the ACM}, 11(6):419--422, 1968.

\bibitem{chomsky56}
N.~Chomsky.
\newblock Three models for the description of language.
\newblock \emph{IRE Transactions on Information Theory}, 2(3):113--124, 1956.

\bibitem{chomsky59}
N.~Chomsky.
\newblock On certain formal properties of grammars.
\newblock \emph{Information and Control}, 2(2):137--167, 1959.

\bibitem{hmu}
J.~E. Hopcroft, R.~Motwani y J.~D. Ullman.
\newblock \emph{Introduction to Automata Theory, Languages, and Computation}.
\newblock 3.\textsuperscript{a} ed. Addison-Wesley, 2006.

\bibitem{sipser}
M.~Sipser.
\newblock \emph{Introduction to the Theory of Computation}.
\newblock 3.\textsuperscript{a} ed. Cengage Learning, 2013.

\bibitem{dragon}
A.~V. Aho, M.~S. Lam, R.~Sethi y J.~D. Ullman.
\newblock \emph{Compilers: Principles, Techniques, and Tools}.
\newblock 2.\textsuperscript{a} ed. Addison-Wesley, 2006.

\end{thebibliography}

\end{document}
